% ===========================================================================
%  Chapitre 15 — Similitudes planes et écriture complexe
%  PE 2026, composante GÉOMÉTRIE
% ===========================================================================
\bmtchapitre{Similitudes planes et écriture complexe}
  {Reconnaître et utiliser une similitude plane directe ou indirecte, et
   passer de l'écriture complexe d'une transformation à ses éléments
   géométriques caractéristiques.}
  {Géométrie \textbullet\ environ 8 heures}

Une isométrie conserve les distances ; une similitude les multiplie toutes par
un même nombre. C'est la transformation qui agrandit ou réduit sans déformer
— celle qui fait passer d'un plan à sa maquette, d'une photographie à son
agrandissement.

Le chapitre tient sa force d'un fait remarquable : dans le plan complexe,
toute similitude directe s'écrit $z' = az+b$, et toute similitude indirecte
$z' = a\overline z+b$. Deux nombres suffisent donc à décrire complètement une
transformation, et l'on passe de la géométrie au calcul par une simple
lecture. C'est ce va-et-vient que le problème de géométrie du baccalauréat
demande chaque année, dans ses parties A et B.

% ---------------------------------------------------------------------------
\begin{revision}
Chapitres 12 et 14.

\begin{enumerate}[label=\textbf{\arabic*.}, leftmargin=*, itemsep=0.35em]
  \item Donner module et argument de $1-i$, de $2i$ et de $-3$.
  \item Rappeler l'interprétation géométrique de
        $\arg\!\left(\frac{z_{C}-z_{A}}{z_{B}-z_{A}}\right)$.
  \item Qu'appelle-t-on une homothétie de centre $\Omega$ et de rapport $k$ ?
  \item Quelle est la nature de la composée d'une rotation et d'une
        homothétie de même centre ?
  \item Résoudre $z = (1-i)z+3$ dans $\C$.
\end{enumerate}
\end{revision}

\begin{automatismes}
Sans calculatrice, sans brouillon.

\begin{enumerate}[label=\textbf{\alph*)}, leftmargin=*, itemsep=0.25em]
  \item Module et argument de $i$
  \item Module et argument de $1+i$
  \item $\abs{-2}$ et $\arg(-2)$
  \item Écriture complexe de la translation de vecteur d'affixe $b$
  \item Écriture complexe de $\homo{O}{3}$
  \item Écriture complexe de $\rot{O}{\pi/2}$
  \item Écriture complexe de la symétrie d'axe $(Ox)$
  \item Rapport de la similitude $z' = 3iz+1$
\end{enumerate}
\end{automatismes}

\begin{corrige}[automatismes]
a) $1$ et $\frac{\pi}{2}$ \quad b) $\sqrt2$ et $\frac{\pi}{4}$ \quad
c) $2$ et $\pi$ \quad d) $z' = z+b$ \quad e) $z' = 3z$ \quad
f) $z' = iz$ \quad g) $z' = \overline z$ \quad h) $3$.
\end{corrige}

% ===========================================================================
\section{Similitude plane directe}

\begin{definition}[similitude]
Une \textbf{similitude} de rapport $k>0$ est une transformation du plan qui
multiplie toutes les distances par $k$ :
$f(M)f(N) = k\,MN$. Elle est \textbf{directe} si elle conserve les angles
orientés, \textbf{indirecte} si elle les change en leurs opposés.
\end{definition}

\begin{remarque}
Une similitude de rapport $1$ est une isométrie : les similitudes prolongent
donc le chapitre précédent. Réciproquement, une similitude directe de rapport
$k \neq 1$ possède toujours un point invariant unique, appelé son
\textbf{centre}.
\end{remarque}

\begin{theoreme}[forme réduite]
Toute similitude directe de rapport $k \neq 1$ est la composée, dans n'importe
quel ordre, d'une homothétie de rapport $k$ et d'une rotation d'angle
$\theta$, toutes deux de centre le centre $\Omega$ de la similitude :
\[
  f = \homo{\Omega}{k} \circ \rot{\Omega}{\theta}
    = \rot{\Omega}{\theta} \circ \homo{\Omega}{k}.
\]
Le triplet $(\Omega,k,\theta)$ constitue ses \textbf{éléments
caractéristiques}.
\end{theoreme}

\begin{methode}{déterminer une similitude directe par deux couples de points}
Une similitude directe est entièrement déterminée par les images de deux
points distincts. Si $f(A) = A'$ et $f(B) = B'$, alors
\[
  k = \frac{A'B'}{AB}, \qquad
  \theta = \left(\vect{AB},\vect{A'B'}\right),
\]
et le centre s'obtient soit par le calcul complexe, soit géométriquement comme
second point d'intersection de deux cercles bien choisis.
\end{methode}

\begin{visualisation}[tourner et agrandir]
\begin{center}
\begin{tikzpicture}[x=1.1cm, y=1.1cm, font=\footnotesize]
  \coordinate (W) at (0,0);
  \coordinate (M) at (2.2,0);
  \coordinate (M1) at (40:2.2);
  \coordinate (M2) at (40:3.4);
  \filldraw[bmtEncre] (W) circle (2pt) node[anchor=north east] {$\Omega$};
  \filldraw[bmtIndigo] (M) circle (2pt) node[anchor=north] {$M$};
  \filldraw[bmtGris] (M1) circle (1.7pt);
  \node[bmtGris, anchor=south east] at (M1) {$M_{1}$};
  \filldraw[bmtLaterite] (M2) circle (2.1pt) node[anchor=south west] {$M'$};
  \draw[bmtGris] (W) -- (M);
  \draw[bmtGris] (W) -- (M2);
  \draw[bmtVetiver, line width=0.9pt, -{Stealth[length=2.4mm]}]
    (0:1.25) arc (0:40:1.25);
  \node[bmtVetiver, anchor=west] at (20:1.4) {$\theta$};
  \draw[bmtLaterite, line width=1.1pt, {Stealth[length=2.2mm]}-{Stealth[length=2.2mm]}]
    (40:2.32) -- (40:3.28);
  \node[bmtLaterite, anchor=west, rotate=40] at (40:2.8) {$\times k$};
\end{tikzpicture}
\end{center}

Une similitude directe se lit en deux temps depuis son centre : on tourne
d'un angle $\theta$, puis on éloigne ou l'on rapproche dans le rapport $k$.
L'ordre n'a aucune importance, puisque rotation et homothétie de même centre
commutent. Retenez cette image : elle rend immédiate la lecture des éléments
caractéristiques sur une figure.
\end{visualisation}

% ===========================================================================
\section{Similitude plane indirecte}

\begin{theoreme}[forme réduite d'une similitude indirecte]
Toute similitude indirecte de rapport $k$ possédant un point invariant
$\Omega$ s'écrit
\[
  f = \homo{\Omega}{k} \circ s_{(\Delta)} = s_{(\Delta)} \circ \homo{\Omega}{k},
\]
où $(\Delta)$ est une droite passant par $\Omega$, appelée \textbf{axe} de la
similitude. Ses éléments caractéristiques sont le centre, le rapport et
l'axe.
\end{theoreme}

\begin{avertissement}
Une similitude indirecte n'a pas d'angle : elle a un axe. Parler de « l'angle
d'une similitude indirecte » est une faute de vocabulaire fréquente, et le
correcteur la relève.
\end{avertissement}

\begin{visualisation}[refléter et agrandir]
\begin{center}
\begin{tikzpicture}[x=1.1cm, y=1.1cm, font=\footnotesize]
  \coordinate (W) at (0,0);
  \draw[bmtGris, line width=0.8pt] (15:-2.6) -- (15:2.6)
    node[anchor=south west]{$(\Delta)$};
  \coordinate (M)  at (70:2.2);
  \coordinate (M1) at (-40:1.9);
  \coordinate (M2) at (-40:3.7);
  \filldraw[bmtEncre] (W) circle (2pt) node[anchor=north east] {$\Omega$};
  \filldraw[bmtIndigo] (M) circle (2pt) node[anchor=south west] {$M$};
  \filldraw[bmtGris] (M1) circle (1.7pt);
  \node[bmtGris] at ($(M1)+(-0.15,-0.35)$) {$M_{1}$};
  \filldraw[bmtLaterite] (M2) circle (2.1pt) node[anchor=north west] {$M'$};
  \draw[bmtGris, dashed] (M) -- (M1);
  \draw[bmtGris] (W) -- (M2);
  \draw[bmtLaterite, line width=1.1pt,
        {Stealth[length=2.2mm]}-{Stealth[length=2.2mm]}]
    (-40:2.15) -- (-40:3.45);
  \node[bmtLaterite] at ($(-40:2.8)+(0.4,0.15)$) {$\times k$};
\end{tikzpicture}
\end{center}
Une similitude indirecte se lit, elle aussi, en deux temps depuis son
centre : on réfléchit par rapport à l'axe $(\Delta)$, puis on éloigne ou l'on
rapproche dans le rapport $k$. À la différence du cas direct, il n'y a pas
d'angle à mesurer : l'axe remplace entièrement l'angle dans la description
du mouvement.
\end{visualisation}

% ===========================================================================
\section{Expression complexe des transformations}

\begin{theoreme}[écriture complexe]
Le plan étant rapporté à un repère orthonormé direct, une transformation $f$
est
\begin{itemize}[leftmargin=*, itemsep=0.1em]
  \item une \textbf{similitude directe} si et seulement si son écriture
        complexe est $z' = az+b$ avec $a \neq 0$ ;
  \item une \textbf{similitude indirecte} si et seulement si son écriture
        complexe est $z' = a\overline z+b$ avec $a \neq 0$.
\end{itemize}
\end{theoreme}

\begin{propriete}[le catalogue des écritures complexes]
\begin{center}
\setlength{\tabcolsep}{9pt}
\renewcommand{\arraystretch}{1.3}
\begin{tabular}{@{}ll@{}}
\toprule
\textsf{\footnotesize\bfseries TRANSFORMATION} &
\textsf{\footnotesize\bfseries ÉCRITURE COMPLEXE} \\
\midrule
Translation de vecteur d'affixe $b$ & $z' = z+b$ \\
Homothétie de centre $\omega$, de rapport $k$ &
  $z' = k(z-\omega)+\omega$ \\
Rotation de centre $\omega$, d'angle $\theta$ &
  $z' = e^{i\theta}(z-\omega)+\omega$ \\
Similitude directe $(\omega,k,\theta)$ &
  $z' = ke^{i\theta}(z-\omega)+\omega$ \\
Réflexion d'axe $(Ox)$ & $z' = \overline z$ \\
Similitude indirecte & $z' = a\overline z+b$ \\
\bottomrule
\end{tabular}
\end{center}
\end{propriete}

\begin{exemple}[trouver une écriture complexe]
Cherchons l'écriture complexe de la similitude directe $s$ de centre
$C(1-2i)$ transformant $A(2-i)$ en $B(3-2i)$.

Le rapport complexe vaut
\[
  a = \frac{z_{B}-z_{C}}{z_{A}-z_{C}}
    = \frac{3-2i-1+2i}{2-i-1+2i}
    = \frac{2}{1+i}
    = 1-i .
\]
L'écriture est donc $z' = (1-i)\left(z-1+2i\right)+1-2i$. On en lit
immédiatement le rapport $\abs{1-i} = \sqrt2$ et l'angle
$\arg(1-i) = -\frac{\pi}{4}$.
\end{exemple}

% ===========================================================================
\section{Des éléments géométriques à partir de l'écriture complexe}

\begin{methode}{lire les éléments d'une similitude directe $z' = az+b$}
\begin{enumerate}[label=\textbf{\arabic*.}, leftmargin=*, itemsep=0.15em]
  \item Si $a = 1$ : c'est la translation de vecteur d'affixe $b$ — pas de
        centre.
  \item Si $a \neq 1$ : le centre $\omega$ est l'unique solution de
        $\omega = a\omega+b$, soit $\omega = \dfrac{b}{1-a}$.
  \item Le rapport vaut $\abs a$, l'angle $\arg a$.
  \item Cas particuliers : $\abs a = 1$ donne une rotation ; $a$ réel
        strictement positif donne une homothétie.
\end{enumerate}
\end{methode}

\begin{methode}{lire les éléments d'une similitude indirecte $z' = a\overline z+b$}
Le rapport vaut $\abs a$. Pour le centre, on résout le système obtenu en
séparant parties réelle et imaginaire dans $\omega = a\overline\omega+b$ — ce
qui donne deux équations à deux inconnues réelles. Lorsqu'un centre existe,
l'axe est la droite passant par ce centre et faisant avec l'axe des abscisses
l'angle $\dfrac{\arg a}{2}$.
\end{methode}

\begin{exemple}[une similitude indirecte]
Soit $\bar S$ d'écriture $z' = (1-i)\overline z+2+4i$.

\emph{Rapport.} $\abs{1-i} = \sqrt2$.

\emph{Centre.} On pose $z = x+iy$ ; alors
$(1-i)(x-iy) = (x-y)-i(x+y)$, et l'équation $z' = z$ devient
\[
  \begin{cases} x = x-y+2 \\ y = -x-y+4 \end{cases}
  \iff
  \begin{cases} y = 2 \\ x = 0 \end{cases}
\]
Le centre est le point d'affixe $2i$.

\emph{Axe.} $\arg(1-i) = -\frac{\pi}{4}$, donc l'axe est la droite passant par
le centre et faisant l'angle $-\frac{\pi}{8}$ avec l'axe des abscisses.
\end{exemple}

\begin{entrainement}[écriture complexe]
\exo[\niveau{1}] Déterminer la nature et les éléments de $z' = 2z+1-i$.

\exo[\niveau{1}] Déterminer la nature et les éléments de $z' = iz+3$.

\exo[\niveau{2}] Déterminer l'écriture complexe de la rotation de centre
$1+i$ et d'angle $\frac{\pi}{3}$.

\exo[\niveau{2}] Déterminer la nature et les éléments de
$z' = (1+i)z+2$.

\exo[\niveau{3}] Déterminer la nature et les éléments de
$z' = 2i\overline z+1$.
\end{entrainement}

\begin{corrige}
\textbf{1.} Le coefficient $2$ est réel positif : l'angle est nul et le
rapport vaut $2$. C'est donc une homothétie. Son centre est le point
invariant : $z = 2z+1-i$ donne $z = -1+i$.

\textbf{2.} $\abs i = 1$ et $\arg i = \frac{\pi}{2}$ : c'est une rotation
d'angle $\frac{\pi}{2}$. Son centre vérifie $z = iz+3$, soit
$z = \dfrac{3}{1-i} = \dfrac32+\dfrac32 i$.

\textbf{3.} $z' = e^{i\pi/3}\left(z-(1+i)\right)+1+i$.

\textbf{4.} $\abs{1+i} = \sqrt2$ et $\arg(1+i) = \frac{\pi}{4}$ : similitude
directe de rapport $\sqrt2$ et d'angle $\frac{\pi}{4}$. Son centre vérifie
$z = (1+i)z+2$, soit $-iz = 2$ et $z = 2i$.

\textbf{5.} L'écriture fait intervenir $\overline z$ : la similitude est
indirecte, de rapport $\abs{2i} = 2$. Comme le rapport diffère de $1$, elle
possède un unique point invariant : en posant $z = x+iy$, le système
$x = 2y+1$ et $y = 2x$ donne $z_{0} = -\frac13-\frac23 i$. Dans le repère
centré en $z_{0}$, l'écriture devient $Z' = 2i\overline Z$ : c'est la composée
de l'homothétie de centre $z_{0}$ et de rapport $2$ avec la réflexion d'axe la
droite passant par $z_{0}$ et de direction $\frac{\pi}{4}$, puisque
$2i = 2e^{2i\pi/4}$.
\end{corrige}

% ===========================================================================
\section{Similitudes et configurations du plan}

\begin{propriete}[ce qu'une similitude conserve]
Une similitude conserve l'alignement, le parallélisme, les milieux, le
contact, les angles géométriques, et transforme une droite en droite, un
cercle en cercle. Elle multiplie les longueurs par $k$ et les aires par
$k^{2}$.
\end{propriete}

\begin{methode}{identifier une similitude sur une configuration}
On cherche deux couples de points $(A,A')$ et $(B,B')$ dont on connaît les
positions, puis
\[
  k = \frac{A'B'}{AB},
  \qquad
  \theta = \left(\vect{AB},\vect{A'B'}\right).
\]
Si un point est manifestement invariant, c'est le centre. Sinon, on passe par
l'écriture complexe : c'est presque toujours plus rapide qu'une construction.
\end{methode}

\begin{entrainement}[configurations]
\exo[\niveau{2}] $ABC$ est rectangle isocèle en $A$. Déterminer la similitude
directe de centre $A$ transformant $B$ en $C$.

\exo[\niveau{2}] $ABCD$ est un carré direct. Déterminer la similitude directe
de centre $A$ transformant $B$ en $C$.

\exo[\niveau{3}] $ABC$ est un triangle tel que $AB = 2$ et $AC = 6$ avec
$\left(\vect{AB},\vect{AC}\right) = \frac{\pi}{2}$. Déterminer les éléments de
la similitude directe de centre $A$ transformant $B$ en $C$, puis son écriture
complexe dans un repère bien choisi.
\end{entrainement}

\begin{corrige}
\textbf{6.} Le rapport vaut $\dfrac{AC}{AB} = 1$ et l'angle
$\left(\vect{AB},\vect{AC}\right) = \pm\frac{\pi}{2}$ : la similitude se
réduit à la rotation de centre $A$ et d'angle $\frac{\pi}{2}$ — ou
$-\frac{\pi}{2}$ selon l'orientation du triangle.

\textbf{7.} Dans un carré direct, $AC = AB\sqrt2$ et
$\left(\vect{AB},\vect{AC}\right) = \frac{\pi}{4}$ : c'est la similitude
directe de centre $A$, de rapport $\sqrt2$ et d'angle $\frac{\pi}{4}$.

\textbf{8.} Le rapport vaut $\dfrac{AC}{AB} = 3$ et l'angle $\frac{\pi}{2}$.
En prenant le repère orthonormé direct de centre $A$ dont le premier vecteur
est $\frac12\vect{AB}$, on a $z_{B} = 1$ et $z_{C} = 3i$ : l'écriture complexe
de la similitude est donc $z' = 3iz$.
\end{corrige}

% ===========================================================================
\begin{annales}
Les similitudes ferment le problème de géométrie, et la question est toujours
la même sous des habillages variés : passer de la configuration à l'écriture
complexe, ou l'inverse. Les énoncés qui suivent sont repris de sujets
réellement posés ; leur provenance est indiquée à chaque fois.

\exoann{Baccalauréat, Madagascar, série S, session 2026 — problème 1, partie I}
$ABC$ est direct, rectangle en $A$, $AC = 4$~cm, $BC = 8$~cm ; $O$ est le
milieu de $[AC]$ et $D$ celui de $[BC]$. Soit $S$ la similitude plane directe
transformant $A$ en $B$ et $O$ en $D$. Déterminer les éléments
caractéristiques de $S$.

\exoann{Baccalauréat, Madagascar, série S, session 2026 — problème 1, partie I, suite}
Avec les mêmes données, $R_{A}$ est la rotation de centre $A$ et d'angle
$-\frac{\pi}{3}$, et l'on pose $f = S \circ R_{A}$.
\begin{enumerate}[label=\textbf{\alph*)}, leftmargin=*, itemsep=0.15em]
  \item Montrer que $f$ est une homothétie dont on précisera le rapport $k$.
  \item Déterminer $f(A)$, puis montrer que le centre $\Omega$ de $f$ est le
        barycentre de $\left\{(A\,;2),(B\,;-1)\right\}$.
\end{enumerate}

\exoann{Baccalauréat, Madagascar, série S, session 2026 — problème 1, partie II}
Déterminer la nature et les éléments caractéristiques de la transformation
$\bar S$ d'expression complexe
\[
  z' = (1-i)\overline z+2+4i .
\]

\exoann{Baccalauréat, Madagascar, série S, session 2022 — problème 1}
$ABCD$ est un carré direct de côté $4$~cm, et $S$ la similitude plane directe
qui laisse invariant le point $A$ et transforme $C$ en $B$.
\begin{enumerate}[label=\textbf{\alph*)}, leftmargin=*, itemsep=0.15em]
  \item Déterminer le rapport et l'angle de $S$.
  \item Le plan étant rapporté au repère
        $\left(A\,;\vect{AB},\vect{AD}\right)$, déterminer l'expression
        complexe de $S$, puis retrouver son rapport et son angle.
\end{enumerate}

\exoann{Baccalauréat, Madagascar, série S, session 2023 — problème 1, partie I}
$ABC$ est rectangle et isocèle en $A$, $AB = 4$~cm. On note
$G = \bary{(B\,;1-\sqrt2),(A\,;\sqrt2)}$ et $\bar S$ la similitude plane
indirecte qui laisse invariant $B$ et transforme $A$ en $C$.
\begin{enumerate}[label=\textbf{\alph*)}, leftmargin=*, itemsep=0.15em]
  \item Vérifier qu'il existe une homothétie $h$ de centre $B$ transformant
        $A$ en $G$, dont on précisera le rapport.
  \item En déduire les éléments caractéristiques de $\bar S$.
\end{enumerate}

\exoann{Baccalauréat, Madagascar, série S, session 2023 — problème 1, partie I, suite}
Avec les mêmes données, $(\mathcal C)$ est le cercle de centre $C$ et de rayon
$2$, $(\Gamma)$ l'ensemble des points $M$ tels que
$MA^{2}-MB^{2}-MC^{2} = -12$, et $I$ le point d'intersection de
$(\mathcal C)$ et $(\Gamma)$ situé à l'extérieur du carré $ABDC$. Montrer que
$I$ est le centre de la similitude plane directe $S$ de rapport $\sqrt3$,
d'angle $\frac{\pi}{2}$, qui transforme $C$ en $D$.

\exoann{Baccalauréat blanc, lycée Philibert Tsiranana, Terminale S, 2022-2023 — problème 1, partie A}
Dans le plan orienté, $OAC$ est équilatéral direct avec $OA = a$, et $B$ est
le point tel que $\vect{OB} = 2\vect{OA}$. Soit $S$ la similitude directe de
centre $C$ qui transforme $A$ en $B$.
\begin{enumerate}[label=\textbf{\alph*)}, leftmargin=*, itemsep=0.15em]
  \item Montrer que $BC = \sqrt3\,AB$.
  \item Déterminer le rapport et l'angle de $S$.
\end{enumerate}

\exoann{Baccalauréat, Cameroun, séries D et TI, session 2024 — exercice 1}
Le plan complexe est muni d'un repère orthonormé. On donne $A$, $B$ et $C$
d'affixes respectives $Z_{A} = 2-i$, $Z_{B} = 3-2i$ et $Z_{C} = 1-2i$.
\begin{enumerate}[label=\textbf{\alph*)}, leftmargin=*, itemsep=0.15em]
  \item Déterminer l'expression complexe de la similitude directe $s$ de
        centre $C$ transformant $A$ en $B$.
  \item Préciser les éléments caractéristiques de $s$.
  \item Quelle est l'image de la droite $(AC)$ par $s$ ?
\end{enumerate}

\exoann{Baccalauréat, Cameroun, série C, session 2023 — exercice 2}
Soit $\Omega$ le point d'affixe $1-2i$, et $B$, $C$ les points d'affixes
respectives $-i$ et $3$.
\begin{enumerate}[label=\textbf{\alph*)}, leftmargin=*, itemsep=0.15em]
  \item Déterminer l'écriture complexe de la similitude directe $S$ de centre
        $\Omega$ telle que $S(C) = B$.
  \item En déduire l'angle de $S$.
\end{enumerate}

\exoann{Baccalauréat, Côte d'Ivoire, série D, session 2023 — exercice 4}
On donne $\Omega(1+i)$, $A(1)$ et $B\!\left(\frac32+\frac i2\right)$. Étudier
la similitude directe $S$ de centre $\Omega$ transformant $A$ en $B$ :
déterminer son écriture complexe, puis son rapport et son angle.
\end{annales}

\begin{corrige}[annales]
\textbf{1.} Le rapport vaut $\dfrac{BD}{AO} = \dfrac{4}{2} = 2$, puisque
$BD = \frac{BC}{2} = 4$ et $AO = \frac{AC}{2} = 2$. L'angle est
$\left(\vect{AO},\vect{BD}\right)$, que la figure donne égal à
$-\frac{\pi}{3}$ — le triangle $CAD$ étant équilatéral. Le centre est le point
invariant, obtenu comme second point d'intersection du cercle de diamètre
$[AB]$ avec le cercle passant par $O$ et $D$ ; l'expression complexe de la
partie II le confirme.

\textbf{2. a)} $f$ est composée de deux similitudes directes : c'est une
similitude directe, de rapport $2 \times 1 = 2$ et d'angle
$-\frac{\pi}{3}+\frac{\pi}{3} = 0$. Un angle nul et un rapport différent de
$1$ caractérisent une homothétie, ici de rapport $k = 2$.
\textbf{b)} $R_{A}(A) = A$, donc $f(A) = S(A) = B$. Le centre $\Omega$
vérifie $\vect{\Omega B} = 2\vect{\Omega A}$, c'est-à-dire
$2\vect{\Omega A}-\vect{\Omega B} = \vecu_{0}$ : c'est exactement la relation
qui définit le barycentre de $\left\{(A\,;2),(B\,;-1)\right\}$.

\textbf{3.} Le rapport vaut $\abs{1-i} = \sqrt2$. Pour le centre, on pose
$z = x+iy$ ; l'égalité $z = (1-i)\overline z+2+4i$ donne
$x = x-y+2$ et $y = -x-y+4$, soit $y = 2$ et $x = 0$ : le centre a pour
affixe $2i$. L'axe passe par ce centre et fait avec l'axe des abscisses
l'angle $\frac{\arg(1-i)}{2} = -\frac{\pi}{8}$. C'est une similitude
indirecte de centre $2i$, de rapport $\sqrt2$ et d'axe cette droite.

\textbf{4. a)} $S$ fixe $A$, donc son rapport vaut
$\frac{AB}{AC} = \frac{4}{4\sqrt2} = \frac{1}{\sqrt2}$, et son angle
$\left(\vect{AC},\vect{AB}\right) = -\frac{\pi}{4}$.
\textbf{b)} Dans le repère $\left(A\,;\vect{AB},\vect{AD}\right)$ on a
$z_{A} = 0$, $z_{B} = 1$, $z_{C} = 1+i$. L'écriture est $z' = az$ avec
$a = \frac{z_{B}}{z_{C}} = \frac{1}{1+i} = \frac{1-i}{2}$. On retrouve
$\abs a = \frac{\sqrt2}{2}$ et $\arg a = -\frac{\pi}{4}$.

\textbf{5. a)} Avec $A(0\,;0)$, $B(4\,;0)$ et $C(0\,;4)$, la somme des
coefficients de $G$ vaut $1$ et $\vect{AG} = (1-\sqrt2)\vect{AB}$, donc
$G\left(4-4\sqrt2\,;0\right)$. Alors
$\frac{\overline{BG}}{\overline{BA}} = \frac{-4\sqrt2}{-4} = \sqrt2$ :
l'homothétie de centre $B$ et de rapport $\sqrt2$ envoie $A$ sur $G$.
\textbf{b)} $\bar S$ fixe $B$ et son rapport vaut
$\frac{BC}{BA} = \frac{4\sqrt2}{4} = \sqrt2$. En écrivant
$\bar S = s \circ h$, la réflexion $s$ doit envoyer $G$ sur $C$ et fixer $B$ :
comme $BG = BC = 4\sqrt2$, son axe est la médiatrice de $[GC]$, qui passe bien
par $B$. Donc $\bar S$ est la similitude indirecte de centre $B$, de rapport
$\sqrt2$ et d'axe la médiatrice de $[GC]$.

\textbf{6.} $I$ appartient à $(\mathcal C)$, donc $IC = 2$ ; il appartient à
$(\Gamma)$, cercle de centre $D$ et de rayon $2\sqrt3$, donc $ID = 2\sqrt3$.
Ainsi $\frac{ID}{IC} = \sqrt3$, ce qui est le rapport annoncé. La position de
$I$ à l'extérieur du carré assure de plus que
$\left(\vect{IC},\vect{ID}\right) = \frac{\pi}{2}$ : le point $I$ est bien le
centre de la similitude directe de rapport $\sqrt3$ et d'angle
$\frac{\pi}{2}$ qui transforme $C$ en $D$.

\textbf{7. a)} $OAC$ est équilatéral de côté $a$ et $\vect{OB} = 2\vect{OA}$,
donc $AB = a$ et $OB = 2a$. Le triangle $OBC$ a pour côtés $OB = 2a$,
$OC = a$ et un angle de $\frac{\pi}{3}$ entre eux ; la loi des cosinus donne
$BC^{2} = 4a^{2}+a^{2}-2 \times 2a \times a \times \frac12 = 3a^{2}$, soit
$BC = a\sqrt3 = \sqrt3\,AB$.
\textbf{b)} $S$ fixe $C$ et envoie $A$ sur $B$ : son rapport vaut
$\frac{CB}{CA} = \frac{a\sqrt3}{a} = \sqrt3$, et son angle
$\left(\vect{CA},\vect{CB}\right)$, que le calcul précédent situe à
$-\frac{\pi}{6}$.

\textbf{8. a)} $a = \dfrac{Z_{B}-Z_{C}}{Z_{A}-Z_{C}} = \dfrac{2}{1+i} = 1-i$,
donc $z' = (1-i)(z-1+2i)+1-2i$.
\textbf{b)} Centre $C(1-2i)$, rapport $\abs{1-i} = \sqrt2$, angle
$-\frac{\pi}{4}$.
\textbf{c)} Une similitude transforme une droite en droite. Comme
$s(C) = C$ et $s(A) = B$, l'image de $(AC)$ est la droite $(BC)$.

\textbf{9. a)} $a = \dfrac{z_{B}-z_{\Omega}}{z_{C}-z_{\Omega}}
= \dfrac{-i-1+2i}{3-1+2i} = \dfrac{-1+i}{2+2i}
= \dfrac{(-1+i)(2-2i)}{8} = \dfrac{4i}{8} = \dfrac i2$. L'écriture complexe
est donc $z' = \frac i2\left(z-1+2i\right)+1-2i$.
\textbf{b)} $\arg\!\left(\frac i2\right) = \frac{\pi}{2}$ : l'angle de $S$ est
$\frac{\pi}{2}$, et son rapport $\frac12$.

\textbf{10.} $a = \dfrac{z_{B}-z_{\Omega}}{z_{A}-z_{\Omega}}
= \dfrac{\frac32+\frac i2-1-i}{1-1-i}
= \dfrac{\frac12-\frac i2}{-i} = \dfrac{1+i}{2}$. L'écriture complexe est
$z' = \frac{1+i}{2}\left(z-1-i\right)+1+i$, le rapport vaut
$\left|\frac{1+i}{2}\right| = \frac{\sqrt2}{2}$ et l'angle
$\frac{\pi}{4}$.
\end{corrige}

% ===========================================================================
\begin{jecherche}[l'agrandissement d'un plan de zébu]
Un artisan grave sur cuivre le plan d'un enclos triangulaire. Sur son
gabarit, l'enclos est le triangle $ABC$ d'affixes $z_{A} = 0$, $z_{B} = 4$ et
$z_{C} = 2+3i$, l'unité valant un centimètre. Il veut en graver une version
agrandie, tournée d'un quart de tour dans le sens direct, de sorte que le
sommet $A$ reste en place et que le côté $[AB]$ mesure $10$~cm.

\begin{enumerate}[label=\textbf{\arabic*.}, leftmargin=*, itemsep=0.3em]
  \item Déterminer le rapport et l'angle de la similitude $S$ cherchée.
  \item Écrire son expression complexe.
  \item Calculer les affixes des images $B'$ et $C'$.
  \item Vérifier que $A B' C'$ a la même forme que $ABC$ : comparer les
        rapports de longueurs.
  \item Quelle est l'aire de l'enclos agrandi, sachant que celle du gabarit
        vaut $6$~cm$^{2}$ ?
  \item L'artisan veut ensuite obtenir une figure retournée, image de $ABC$
        par une similitude indirecte de centre $A$, de rapport $\frac52$, et
        d'axe l'axe des abscisses. Donner son expression complexe et l'affixe
        de l'image de $C$.
\end{enumerate}
\end{jecherche}

\begin{corrige}[l'agrandissement d'un plan de zébu]
\textbf{1.} $AB = 4$ doit devenir $10$ : le rapport vaut
$k = \frac{10}{4} = \frac52$. L'angle est $\frac{\pi}{2}$, et le centre est
$A$, d'affixe $0$.

\textbf{2.} $z' = \frac52\,e^{i\pi/2}\,z = \frac52\,i\,z$.

\textbf{3.} $z_{B'} = \frac52 i \times 4 = 10i$ et
$z_{C'} = \frac52 i(2+3i) = -\frac{15}{2}+5i$.

\textbf{4.} Une similitude multiplie toutes les longueurs par le même
rapport : $\frac{AB'}{AB} = \frac{AC'}{AC} = \frac{B'C'}{BC} = \frac52$. Les
rapports internes sont donc inchangés, et les deux triangles ont la même
forme.

\textbf{5.} Les aires sont multipliées par $k^{2} = \frac{25}{4}$ :
l'enclos agrandi a pour aire $6 \times \frac{25}{4} = 37{,}5$~cm$^{2}$.

\textbf{6.} Une similitude indirecte de centre $A$ — d'affixe nulle —, de
rapport $\frac52$ et d'axe l'axe des abscisses a pour expression
$z' = \frac52\overline z$. L'image de $C$ a donc pour affixe
$\frac52(2-3i) = 5-\frac{15}{2}i$.
\end{corrige}

% ===========================================================================
\begin{jemeteste}
Répondre par vrai ou faux, en justifiant en une ligne.

\begin{enumerate}[label=\textbf{\arabic*.}, leftmargin=*, itemsep=0.28em]
  \item Toute similitude directe possède un centre.
  \item Une similitude de rapport $1$ est une isométrie.
  \item L'écriture $z' = 3\overline z+i$ définit une similitude directe.
  \item Une similitude multiplie les aires par son rapport.
  \item Une similitude indirecte possède un angle.
  \item Une similitude transforme un cercle en cercle.
\end{enumerate}
\end{jemeteste}

\begin{corrige}[je me teste]
\textbf{1.} Faux — les translations sont des similitudes directes sans
centre.
\textbf{2.} Vrai, par définition.
\textbf{3.} Faux — la présence du conjugué en fait une similitude indirecte.
\textbf{4.} Faux — elle les multiplie par le carré du rapport.
\textbf{5.} Faux — elle a un axe, pas un angle.
\textbf{6.} Vrai — le centre est l'image du centre, le rayon est multiplié
par le rapport.
\end{corrige}

% ===========================================================================
\begin{commeaubac}
\bmtsource{Baccalauréat, série S, session 2026 — problème 1}
$ABC$ est un triangle direct rectangle en $A$ tel que $AC = 4$~cm et
$BC = 8$~cm ; $O$ est le milieu de $[AC]$ et $D$ celui de $[BC]$. Soit $S$ la
similitude plane directe qui transforme $A$ en $B$ et $O$ en $D$, et $R_{A}$
la rotation de centre $A$ et d'angle $-\frac{\pi}{3}$. On pose
$f = S \circ R_{A}$.

\begin{enumerate}[label=\textbf{\arabic*.}, leftmargin=*, itemsep=0.25em]
  \item Déterminer les éléments caractéristiques de $S$.
        \hfill\textup{(0,75 pt)}
  \item Montrer que $f$ est une homothétie de rapport $k$ à préciser.
        \hfill\textup{(0,5 pt)}
  \item Déterminer $f(A)$ et en déduire que le centre $\Omega$ de $f$ est le
        barycentre de $\left\{(A\,;2),(B\,;-1)\right\}$.
        \hfill\textup{(0,5 pt)}
  \item Le plan étant rapporté à un repère complexe convenable, déterminer la
        nature et les éléments de la transformation d'expression
        $z' = (1-i)\overline z+2+4i$. \hfill\textup{(0,75 pt)}
\end{enumerate}
\end{commeaubac}

\begin{corrige}[comme au baccalauréat]
\textbf{1.} $S$ envoie $A$ sur $B$ et $O$ sur $D$ : son rapport vaut
$\frac{BD}{AO} = \frac{4}{2} = 2$ et son angle
$\left(\vect{AO},\vect{BD}\right) = -\frac{\pi}{3}$, le triangle $CAD$ étant
équilatéral.

\textbf{2.} $f$ est une similitude directe de rapport $2 \times 1 = 2$ et
d'angle $-\frac{\pi}{3}+\frac{\pi}{3} = 0$ : c'est une homothétie de rapport
$k = 2$.

\textbf{3.} $f(A) = S\left(R_{A}(A)\right) = S(A) = B$. Le centre vérifie
donc $\vect{\Omega B} = 2\vect{\Omega A}$, soit
$2\vect{\Omega A}-\vect{\Omega B} = \vecu_{0}$ : $\Omega$ est le barycentre
de $\left\{(A\,;2),(B\,;-1)\right\}$, c'est-à-dire le symétrique de $B$ par
rapport à $A$.

\textbf{4.} Le conjugué signale une similitude indirecte, de rapport
$\abs{1-i} = \sqrt2$. Le centre, solution de $z = (1-i)\overline z+2+4i$, a
pour affixe $2i$ ; l'axe passe par ce point et fait l'angle
$-\frac{\pi}{8}$ avec l'axe des abscisses.
\end{corrige}

% ===========================================================================
\begin{concours}
\bmtsource{D'après un concours d'entrée en école d'ingénieurs}
On considère la transformation $f$ d'écriture complexe
$z' = az+b$, avec $a \neq 1$, et l'on note $\Omega$ son centre.

\begin{enumerate}[label=\textbf{\arabic*.}, leftmargin=*, itemsep=0.25em]
  \item Montrer que pour tout point $M$ d'affixe $z$,
        $z'-\omega = a(z-\omega)$.
  \item En déduire que $f^{n}$, composée de $f$ avec elle-même $n$ fois, est
        la similitude directe de centre $\Omega$, de rapport $\abs a^{n}$ et
        d'angle $n\arg a$.
  \item On prend $a = \frac{1+i}{2}$ et $\Omega$ l'origine. Étudier la suite
        des points $M_{n}$ définie par $M_{0}$ d'affixe $1$ et
        $M_{n+1} = f(M_{n})$ : décrire son comportement.
\end{enumerate}
\end{concours}

\begin{corrige}[vers le concours]
\textbf{1.} Le centre vérifie $\omega = a\omega+b$. En soustrayant cette
égalité de $z' = az+b$, il vient $z'-\omega = a(z-\omega)$.

\textbf{2.} Une récurrence immédiate donne
$f^{n}(z)-\omega = a^{n}(z-\omega)$ : c'est l'écriture d'une similitude
directe de centre $\Omega$, de rapport $\abs{a^{n}} = \abs a^{n}$ et d'angle
$\arg\!\left(a^{n}\right) = n\arg a$.

\textbf{3.} Ici $\abs a = \frac{\sqrt2}{2}<1$ et $\arg a = \frac{\pi}{4}$.
L'affixe de $M_{n}$ vaut $a^{n}$ : le point tourne d'un huitième de tour à
chaque étape tout en se rapprochant de l'origine dans le rapport
$\frac{\sqrt2}{2}$. La suite des points décrit une spirale qui converge vers
$\Omega$, et l'on repasse par la même direction tous les huit termes.
\end{corrige}

% ===========================================================================
\begin{notepourlaclasse}
Le chapitre est court et rentable, à condition d'ancrer une seule idée : dans
le plan complexe, une similitude directe est une multiplication. Le rapport
est un module, l'angle un argument, le centre un point fixe. Les élèves qui
tiennent cette phrase traitent la partie B du problème en dix minutes.

L'erreur de vocabulaire la plus fréquente consiste à attribuer un angle à une
similitude indirecte. Corrigez-la dès la première occurrence : une similitude
indirecte a un axe, et cet axe fait avec l'axe des abscisses la moitié de
l'argument de $a$ — résultat qu'il vaut mieux faire redémontrer une fois que
faire apprendre.

Enfin, entraînez au passage inverse : partir de la configuration, choisir un
repère, écrire les affixes, en déduire $a$ et $b$. C'est ce sens-là qui est
demandé en début de partie B, et c'est celui que les élèves travaillent le
moins.
\end{notepourlaclasse}

% =========================================================================
\begin{devoir}[2 heures]
Trois exercices indépendants, notés sur $20$. Le devoir est cumulatif : il porte sur ce chapitre et sur ceux qui le précèdent. Les énoncés sont repris de sessions réelles et assemblés ici en épreuve ; leur provenance est indiquée à chaque fois.

\exods{1}{Baccalauréat, Madagascar, série S, session 2022 — problème 1, partie A}{5 points}
$ABCD$ est un carré direct de côté $4$~cm. On note $t$ la translation de
vecteur $\vect{AC}$, $t_{1}$ celle de vecteur $\vect{DC}$, $t_{2}$ celle de
vecteur $\vect{AD}$, $s_{(AD)}$ la réflexion d'axe $(AD)$, et
$f = t \circ s_{(AD)}$.
\begin{enumerate}[label=\textbf{\alph*)}, leftmargin=*, itemsep=0.15em]
  \item Justifier que $t = t_{1} \circ t_{2}$.
  \item Décomposer $t_{1}$ en deux symétries orthogonales dont l'un des axes
        est la droite $(AD)$.
  \item En déduire la nature et les éléments caractéristiques de $f$.
\end{enumerate}

\exods{2}{Baccalauréat blanc, lycée Philibert Tsiranana, Terminale S, 2022-2023 — problème 1, partie A}{7 points}
$OAC$ est un triangle équilatéral tel que
$\left(\vect{OA},\vect{OC}\right) = \frac{\pi}{3}$ et $OA = a$. On note
$R_{O}$ la rotation de centre $O$ et d'angle $\frac{\pi}{3}$, $R_{A}$ celle de
centre $A$ et d'angle $\frac{2\pi}{3}$, et $R = R_{O} \circ R_{A}$.
\begin{enumerate}[label=\textbf{\alph*)}, leftmargin=*, itemsep=0.15em]
  \item Montrer que $R$ est une symétrie centrale.
  \item Déterminer $R(A)$ et en déduire le centre de $R$.
  \item Déterminer les droites $(D_{1})$ et $(D_{2})$ telles que
        $R_{O} = s_{(D_{2})} \circ s_{(OA)}$ et
        $R_{A} = s_{(OA)} \circ s_{(D_{1})}$, puis retrouver les résultats
        précédents.
\end{enumerate}

\exods{3}{Baccalauréat, Madagascar, série S, session 2022 — problème 1}{8 points}
$ABCD$ est un carré direct de côté $4$~cm, et $S$ la similitude plane directe
qui laisse invariant le point $A$ et transforme $C$ en $B$.
\begin{enumerate}[label=\textbf{\alph*)}, leftmargin=*, itemsep=0.15em]
  \item Déterminer le rapport et l'angle de $S$.
  \item Le plan étant rapporté au repère
        $\left(A\,;\vect{AB},\vect{AD}\right)$, déterminer l'expression
        complexe de $S$, puis retrouver son rapport et son angle.
\end{enumerate}

\end{devoir}

\begin{corrige}[devoir surveillé]
\textbf{Exercice 1. a)} $\vect{AC} = \vect{AD}+\vect{DC}$, donc
$t = t_{1} \circ t_{2}$, les translations commutant.
\textbf{b)} $t_{1}$, de vecteur $\vect{DC}$ perpendiculaire à $(AD)$, se
décompose en $s_{(\Delta)} \circ s_{(AD)}$ où $(\Delta)$ est parallèle à
$(AD)$, à la distance $\frac{DC}{2} = 2$.
\textbf{c)} Alors
$f = t_{1} \circ t_{2} \circ s_{(AD)}
= s_{(\Delta)} \circ s_{(AD)} \circ t_{2} \circ s_{(AD)}$. Comme $t_{2}$ a
pour vecteur $\vect{AD}$, dirigé par l'axe, la composée
$t_{2} \circ s_{(AD)}$ est une symétrie glissée d'axe $(AD)$ et de vecteur
$\vect{AD}$. Au total, $f$ est une symétrie glissée d'axe $(\Delta)$ et de
vecteur $\vect{AD}$.

\textbf{Exercice 2. a)} La somme des angles vaut $\frac{\pi}{3}+\frac{2\pi}{3} = \pi$ :
$R$ est une symétrie centrale.
\textbf{b)} $R(A) = R_{O}\left(R_{A}(A)\right) = R_{O}(A)$, image de $A$ par
la rotation de centre $O$ et d'angle $\frac{\pi}{3}$ : c'est $C$. Le centre de
la symétrie centrale est donc le milieu de $[AC]$.
\textbf{c)} $(D_{2})$ est l'image de $(OA)$ par la rotation de centre $O$ et
d'angle $\frac{\pi}{6}$ ; $(D_{1})$ est la droite passant par $A$ formant
l'angle $-\frac{\pi}{3}$ avec $(OA)$. La composée redonne bien la symétrie
centrale de centre le milieu de $[AC]$.

\textbf{Exercice 3. a)} $S$ fixe $A$, donc son rapport vaut
$\frac{AB}{AC} = \frac{4}{4\sqrt2} = \frac{1}{\sqrt2}$, et son angle
$\left(\vect{AC},\vect{AB}\right) = -\frac{\pi}{4}$.
\textbf{b)} Dans le repère $\left(A\,;\vect{AB},\vect{AD}\right)$ on a
$z_{A} = 0$, $z_{B} = 1$, $z_{C} = 1+i$. L'écriture est $z' = az$ avec
$a = \frac{z_{B}}{z_{C}} = \frac{1}{1+i} = \frac{1-i}{2}$. On retrouve
$\abs a = \frac{\sqrt2}{2}$ et $\arg a = -\frac{\pi}{4}$.

\end{corrige}
